% Prospective support draft, 2026-10-18. NOT a completed experiment report.
% Not integrated into any immutable manuscript version; not covered by v6 audit.
% Governing protocol: #1185 amendments 7 (revision 1 + adoption corrections), 8.
% Amendment 8 text: ca25a11bda465add9240ac794ed787a081d50772,
% refs/amendment8-physical-baseline.md.
% Final selection: comments 6247383821 + 6247421487.
% Start sampling: comment 6247585770 + subsequent rationale correction.
% Scope decisions: Squad reviews 18, 19, 21; EXPERIMENTS.md.
% Before incorporation: verify against merged implementation and main manifests;
% insert real host/environment and results; resolve pending markers; new audit.
% Render as a fragment inside article with amsmath and booktabs, or the paper class.

\section{Prospective optimization study}
\label{sec:prospective-optimization}

\noindent\textbf{Draft status.}
This section describes the adopted protocol for an experiment that has not yet
been run. It reports no optimization outcome or speedup. The final manuscript
will need the measured results, execution metadata and independent refinement
checks before making a performance claim.

\subsection{Question and comparison}

The study will test whether exact discrete field gradients reduce the cost of
reaching a common validated torque-ripple target, and whether that result
changes when additional shape freedom is available. A coarse nonlinear FEM
model provides the inexpensive physical competitor. Comparing both optimizers
at both training resolutions separates optimizer choice from spatial
resolution within this experiment (Table~\ref{tab:planned-methods}). The
optimizer comparison includes differences in search strategy; it does not
isolate gradient access as the only causal difference.

\begin{table}[ht]
\centering
\begin{tabular}{lll}
\toprule
Method & Optimizer or search model & Training FEM level \\
\midrule
A & Safeguarded SLSQP, field adjoints & L1 \\
B & COBYQA, forward evaluations & L1 \\
C & Empirical cubic-RBF trust-region model & L1 \\
D & COBYQA, forward evaluations & L0 \\
E & Safeguarded SLSQP, field adjoints & L0 \\
\bottomrule
\end{tabular}
\caption{Planned methods. A/B/D/E form the optimizer-by-resolution comparison.
C is an empirical surrogate fitted to field evaluations. D and E solve the
coarse nonlinear field equations; they are neither fitted magnetic circuits
nor reduced-basis models.}
\label{tab:planned-methods}
\end{table}

Both training levels will use harmonic maximum $h=45$, twenty rotor angles,
the same nonlinear material law, and current fractions
$\mathcal C=\{0,\tfrac12,1\}$ of the fixed peak current $9.6777\,\mathrm{A}$.
Only spatial resolution changes between L0 and L1. Each level will have its own
initial-design training normalization, measured before the optimization runs.
There is no fitted correction from coarse to fine objective values.
Method C's post-run current predictions use the separately specified linear
interpolation between its current-specific RBF predictions; these predictions
are not nonlinear physical solves at the withheld current.
For D and E, each stopped run will select its lowest-objective common-validated
feasible point, with the earliest validation breaking ties. If none exists,
the last buildable training incumbent will be used with an explicit lack of
validated feasibility; if that is also missing, the diagnostic is unavailable.
At this frozen point, L0 and the common validation model will solve the
nonlinear equations at fractions $0.25$ and $0.75$ on twenty angles, retaining
waveforms, means and standard deviations. These checks will not fit a correction
or modify the stopped optimization.

\subsection{Design spaces, objective and feasibility}

Space A has eight free engineering parameters. Space B contains the same
scalars and ten grouped surface-control offsets acting on 55 control points;
the points are not independent design variables. Axial length, overall radial
scale, stator diameter and operating-angle convention remain fixed. The
parameter box bounds magnet area but does not hold magnet volume constant.
The continuous-geometry gap and outer-web bounds established separately do
not constitute manufacturing qualification or a guarantee for every mesh.

Let $\sigma_{r,c}(\theta)$ and $\mu_{r,c}(\theta)$ denote torque standard
deviation and mean at design $\theta$, model $r$, and current fraction $c$.
The training objective and full-current constraint are
\begin{align}
 F_r(\theta) &= \sum_{c\in\mathcal C}\sigma_{r,c}(\theta)^2,
 & \mu_{r,1}(\theta)&\geq 1.0025\,\mu_{r,1}(\theta_0).
\end{align}
Optimization uses unit-box coordinates and objective/mean normalization by
the corresponding level's initial-design values. Only the full-current mean
is constrained. Half-current mean torque will be reported as a diagnostic;
this protocol supplies no partial-load mean-torque guarantee.

Construction margins $m_i$ have fixed positive scales $s_i$ of matching units.
The additional normalized constraints are
\begin{equation}
 c_i(\theta)=m_i(\theta)/s_i,\qquad
 h_i(\theta)=c_i(\theta)-0.05\geq0.
\end{equation}
The construction's actual strict and non-strict guards and its discrete
checks determine whether a trial is buildable. Buildable trials can still
violate $h_i\geq0$. The value $0.05$ is a prospective domain margin, not a
certified physical safety factor. An undefined construction or derivative
will be rejected before the corresponding field request.

SLSQP will abort a call that requests a refused point, restore toward the last
admissible evaluated point by at most twenty bisections, and restart at the
actual restored point with fresh optimizer state. At most twenty restarts
are allowed. COBYQA will use the pinned NaN barrier. The empirical surrogate
will skip refused samples and use the registered coordinate-sign retries and
halvings to construct its initial set. Its feasible-first ordering must include
the geometry constraints as well as the mean constraint.
Moved meshes and any fresh replacements must satisfy the common mesh policy.
Fresh repair permits four attempts with the registered, non-cumulative
settings. Every attempt, refusal and restoration will be recorded and charged.

\subsection{Common validation and target credit}

Every method will use the same L2/$h93$/g20 validation model and initial-design
anchor. Write its initial objective as $F_0$ and its full-current mean as
$\mu_0$. For a completed candidate validation, define
\begin{equation}
 \delta F(\theta)=\max\left\{0.10F_{\mathrm{val}}(\theta),
                    2.5\times10^{-7}\,(\mathrm{N\,m})^2\right\}.
\end{equation}
The three fixed targets are $T_{50}=0.5F_0$, $T_{90}=0.1F_0$, and
$T_{99}=0.01F_0$. A target is attained only when
\begin{equation}
 F_{\mathrm{val}}(\theta)+\delta F(\theta)\leq T,
 \qquad \mu_{\mathrm{val},1}(\theta)\geq1.0025\mu_0,
\end{equation}
the geometry/discrete checks and validation-mesh policy pass, and validation
finishes within budget. The allowance is an empirical validation rule, not
a rigorous bound on continuum error. Optimizer success flags confer no credit.
The registered callback rules validate changed incumbents, excluding the
start, duplicate validations and incomplete initial interpolation sets.

Each run will receive 120 measured cold-validation wall-time units, with a
60-unit checkpoint using only validations completed by that checkpoint.
Both training resolutions share the same validation timing unit; L0 costs
will be measured without a nominal resolution discount. A run that later
fails retains its earlier eligible target crossings. Failed, deferred,
contaminated and budget-censored outcomes will remain visible.

\subsection{Starts, scheduling and cost}

The planned cohort is $\theta_0$ plus four seeded starts in each space, five
methods and two timing repeats: 100 optimization runs. Each seed uses the
registered ten-point, eighteen-dimensional Latin-hypercube sequence.
The scalar coordinates are matched across
spaces. The frozen pilot indices are tried first for seeds 1 and 2; deterministic
replacement indices and the selected indices for seeds 3 and 4 will be archived.
A start must pass admissibility/buildability checks and the moved-reference-mesh
policy without remeshing in both spaces at L0 and L1. This favors moderate
deformations reachable by mesh motion. The starts are not uniform samples of
the admissible domain, although subsequent optimizer iterates may use repairs.
The comparison is a paired case study over five designs; timing repeats are
not independent design samples.

For start index $s$, repeat 1 will rotate the method order ABCDE left by $s$;
repeat 2 will reverse that exact order. With concurrent workers, this is the
queue-entry order. The actual dispatch/completion order and worker count will
be retained. The host gate and load monitor apply to the common calibration
and to optimization runs. Post-run diagnostics have a separate interval and
cost record. A run will not launch after a failed idle gate; contaminated runs
will retain their records and be rerun under the adopted policy.

The charged run ledger includes evaluator startup, geometry checks, mesh
updates and repairs, assembly and solves, adjoints, search overhead,
surrogate fitting and validation. Shared preparation/calibration costs and
post-run costs will also be reported, together with end-to-end execution
costs. CPU seconds will be measured with explicit process/child boundaries.
RSS observations will identify their process scope and cumulative high-water
semantics; separate process maxima will not be summed into a purported
simultaneous process-tree peak. [PENDING main-run host, process count,
environment versions, calibration records and source/data hashes]

\subsection{Final checks and interpretation}

For each method and space, the best-observed representative will be the
lowest-$F_{\mathrm{val}}$ eligible feasible validation, breaking ties by start
index, repeat index and validation ordinal. Earlier eligible records of
failed runs remain candidates. A missing feasible representative will be
reported explicitly. For each method, space and target, the earliest attaining
validation across starts in repeat 1 will also be selected: at most thirty
crossing designs in addition to ten representatives. A missing repeat-1
crossing will not be replaced with one from repeat 2.

The initial design and these selected points will be checked at L3/$h93$/g20
and L2/$h93$/g40, retaining the training-current waveforms and the specified
half-current and $0.25/0.75$ diagnostics. Conforming outer-band Arkkio checks
are restricted to fresh meshes with zero stator offsets and supported
rotations. For other designs the missing conforming check and resulting
model dependence will be stated. Exact geometry/settings identities may
share a check while retaining every source role. A failed refinement will
not trigger substitution of a different selected candidate.

Time-to-target statements will refer to common L2/$h93$/g20 validation.
The selected refinement checks will qualify those statements where they
disagree; they do not certify every crossing or typical-start behavior.
Paired timing comparisons will retain censored targets and avoid a finite
speedup ratio when a comparator does not attain the target. Conclusions will
be confined to the tested optimizers, physical models, design spaces and
sampling conditions. Population methods such as CMA-ES remain untested, and
differences between spaces A and B will not be described as a measurement of
intrinsic effective dimension.
